\documentclass[11pt]{article}
\usepackage{amsmath,amssymb,amsthm}
\usepackage{bm}
\usepackage{ctex}
\usepackage{geometry}
\geometry{margin=1in}

\newcommand{\R}{\mathbb{R}}
\newcommand{\C}{\mathbb{C}}
\renewcommand{\top}{\mathrm{T}}
\newcommand{\ii}{\mathrm{i}}
\newcommand{\diag}{\operatorname{diag}}
\newcommand{\rank}{\operatorname{rank}}
\newcommand{\norm}[1]{\left\lVert #1 \right\rVert}
\newcommand{\abs}[1]{\left\lvert #1 \right\rvert}

\newtheorem{lemma}{引理}
\newtheorem{theorem}{定理}
\newenvironment{solution}{\par\noindent\textbf{解：}}{\par}

\begin{document}

\section*{第一次作业解答}

\begin{enumerate}

\item 设
\[
x=(1,0,4,6,3,4)^{\top}.
\]
求一个实 Householder 变换 \(H\in\R^{6\times 6}\) 和一个正数 \(\alpha\)，使得
\[
Hx=(1,\alpha,4,6,0,0)^{\top}.
\]

\begin{solution}
Householder 变换保持向量 \(2\) 范数，因此
\[
\norm{Hx}_2=\norm{x}_2.
\]
而
\[
\norm{x}_2^2=1^2+0^2+4^2+6^2+3^2+4^2=78.
\]
目标向量
\[
y=(1,\alpha,4,6,0,0)^{\top}
\]
的范数平方为
\[
\norm{y}_2^2=1+\alpha^2+16+36=53+\alpha^2.
\]
于是
\[
53+\alpha^2=78,
\]
故
\[
\alpha^2=25.
\]
因为 \(\alpha>0\)，所以
\[
\alpha=5.
\]

下面构造 Householder 变换。令
\[
z=(1,0,4,6)^{\top},\qquad w=(1,5,4,6)^{\top}.
\]
取
\[
v=z-w=(0,-5,0,0)^{\top}.
\]
令
\[
u=\frac{v}{\norm{v}_2}=(0,-1,0,0)^{\top}.
\]
定义
\[
\tilde H=I_4-2uu^{\top}
=
\begin{pmatrix}
1&0&0&0\\
0&-1&0&0\\
0&0&1&0\\
0&0&0&1
\end{pmatrix}.
\]
则
\[
\tilde H z=w.
\]
于是取
\[
H=
\begin{pmatrix}
\tilde H&0\\
0&I_2
\end{pmatrix}
=
\diag(1,-1,1,1,1,1).
\]
验证：
\[
Hx=(1,5,4,6,0,0)^{\top}.
\]
因此
\[
\boxed{\alpha=5,\qquad H=\diag(1,-1,1,1,1,1).}
\]
\end{solution}

\item 设
\[
x=(1,1-2\ii,2+\ii,2+\ii,0)^{\top},
\]
求一个复 Householder 变换 \(H\in\C^{5\times 5}\) 和一个数 \(\alpha\in\C\)，使得
\[
Hx=(\alpha,0,1,2+\ii,0)^{\top}.
\]

\begin{solution}
复 Householder 变换保持复欧几里得范数，因此
\[
\norm{Hx}_2=\norm{x}_2.
\]
计算
\[
\norm{x}_2^2=1+\abs{1-2\ii}^2+\abs{2+\ii}^2+\abs{2+\ii}^2
=1+5+5+5=16.
\]
目标向量
\[
y=(\alpha,0,1,2+\ii,0)^{\top}
\]
满足
\[
\norm{y}_2^2=\abs{\alpha}^2+1+5=\abs{\alpha}^2+6.
\]
于是
\[
\abs{\alpha}^2+6=16,
\]
故
\[
\abs{\alpha}^2=10.
\]
可取
\[
\alpha=\sqrt{10}.
\]

令
\[
z=x,\qquad w=(\sqrt{10},0,1,2+\ii,0)^{\top},
\]
并取
\[
v=z-w.
\]
复 Householder 矩阵为
\[
H=I-2\frac{vv^{H}}{v^{H}v},
\]
其中 \(v^{H}\) 表示共轭转置。由 Householder 变换性质，
\[
Hz=w.
\]
因此
\[
\boxed{
H=I-2\frac{vv^{H}}{v^{H}v},\qquad
v=(1-\sqrt{10},\,1-2\ii,\,1+\ii,\,0,\,0)^{\top},
}
\]
并且
\[
\boxed{\alpha=\sqrt{10}.}
\]
若取相反符号，则 \(\alpha=-\sqrt{10}\)，对应
\[
v=x-(-\sqrt{10},0,1,2+\ii,0)^{\top}.
\]
\end{solution}

\item 确定 \(c=\cos\theta\) 和 \(s=\sin\theta\)，使得
\[
\begin{pmatrix}
c&s\\
-s&c
\end{pmatrix}
\begin{pmatrix}
5\\
12
\end{pmatrix}
=
\alpha
\begin{pmatrix}
1\\
1
\end{pmatrix},
\qquad \alpha\in\R.
\]

\begin{solution}
设
\[
\begin{pmatrix}
c&s\\
-s&c
\end{pmatrix}
\begin{pmatrix}
5\\
12
\end{pmatrix}
=
\begin{pmatrix}
5c+12s\\
-5s+12c
\end{pmatrix}
=
\begin{pmatrix}
\alpha\\
\alpha
\end{pmatrix}.
\]
于是
\[
\begin{cases}
5c+12s=\alpha,\\
-5s+12c=\alpha.
\end{cases}
\]
两式相减得
\[
5c+12s+5s-12c=0,
\]
即
\[
17s-7c=0.
\]
因此
\[
17s=7c.
\]
又因为
\[
c^2+s^2=1,
\]
将 \(s=\frac{7}{17}c\) 代入，得
\[
c^2+\frac{49}{289}c^2=1,
\]
即
\[
\frac{338}{289}c^2=1.
\]
所以
\[
c^2=\frac{289}{338}.
\]
因此
\[
c=\pm \frac{17}{\sqrt{338}},\qquad
s=\pm \frac{7}{\sqrt{338}}.
\]
由于 \(17s=7c\)，\(c\) 与 \(s\) 同号。因此两组解为
\[
\boxed{
c=\frac{17}{\sqrt{338}},\qquad
s=\frac{7}{\sqrt{338}},
}
\]
或
\[
\boxed{
c=-\frac{17}{\sqrt{338}},\qquad
s=-\frac{7}{\sqrt{338}}.
}
\]
对应
\[
\alpha=5c+12s.
\]
若取第一组，则
\[
\alpha=\frac{169}{\sqrt{338}};
\]
若取第二组，则
\[
\alpha=-\frac{169}{\sqrt{338}}.
\]
\end{solution}

\item 设 \(L\in\R^{m\times n}\)（\(m\geqslant n\)）是下三角阵，说明如何确定 Householder 矩阵
\(H_1,H_2,\dots,H_n\in\R^{m\times m}\)，使得
\[
H_n\cdots H_2H_1L=
\begin{pmatrix}
L_1\\
O
\end{pmatrix},
\]
其中 \(L_1\in\R^{n\times n}\) 是下三角矩阵。

\begin{solution}
设
\[
L=
\begin{pmatrix}
l_{11}&0&\cdots&0\\
l_{21}&l_{22}&\cdots&0\\
\vdots&\vdots&\ddots&\vdots\\
l_{m1}&l_{m2}&\cdots&l_{mn}
\end{pmatrix}.
\]
目标是逐列消去第 \(n+1\) 行到第 \(m\) 行的元素，同时保持前 \(n\) 行为下三角形式。

第 \(1\) 步：考察第一列从第 \(2\) 行到第 \(m\) 行的元素
\[
\begin{pmatrix}
l_{21}\\
l_{31}\\
\vdots\\
l_{m1}
\end{pmatrix}.
\]
构造 Householder 矩阵 \(\widetilde H_1\in\R^{(m-1)\times(m-1)}\)，使得
\[
\widetilde H_1
\begin{pmatrix}
l_{21}\\
l_{31}\\
\vdots\\
l_{m1}
\end{pmatrix}
=
\begin{pmatrix}
*\\
0\\
\vdots\\
0
\end{pmatrix}.
\]
令
\[
H_1=
\begin{pmatrix}
1&0\\
0&\widetilde H_1
\end{pmatrix}
\in\R^{m\times m}.
\]
则 \(H_1L\) 的第一列只有前两个元素可能非零，其余为零。

第 \(2\) 步：考察第二列从第 \(3\) 行到第 \(m\) 行的元素。构造
\(\widetilde H_2\in\R^{(m-2)\times(m-2)}\)，使得这些元素被消去。令
\[
H_2=
\begin{pmatrix}
I_2&0\\
0&\widetilde H_2
\end{pmatrix}
\in\R^{m\times m}.
\]
由于 \(H_2\) 只作用在第 \(3\) 到第 \(m\) 行，不会破坏第一列已经消好的零，也不会破坏下三角结构。

依此类推。第 \(k\) 步消去第 \(k\) 列中第 \(k+1\) 到第 \(m\) 行的元素，构造
\[
H_k=
\begin{pmatrix}
I_k&0\\
0&\widetilde H_k
\end{pmatrix},
\qquad
\widetilde H_k\in\R^{(m-k)\times(m-k)}.
\]
经过 \(n\) 步后得到
\[
H_nH_{n-1}\cdots H_1L=
\begin{pmatrix}
L_1\\
O
\end{pmatrix},
\]
其中 \(L_1\in\R^{n\times n}\) 是下三角矩阵。

因此算法为：
\[
\boxed{
\text{对 }k=1,2,\dots,n,\text{ 用 Householder 变换消去第 }k\text{ 列第 }k+1,\dots,m\text{ 行元素。}
}
\]
\end{solution}

\item 设 \(H\in\C^{n\times n}\) 是上 Hessenberg 矩阵，\(R\in\C^{n\times n}\) 是上三角矩阵。证明：\(HR\) 是上 Hessenberg 矩阵。

\begin{solution}
上 Hessenberg 矩阵定义为
\[
H_{ij}=0,\qquad i>j+1.
\]
上三角矩阵满足
\[
R_{ij}=0,\qquad i>j.
\]
记
\[
C=HR.
\]
需要证明
\[
C_{ij}=0,\qquad i>j+1.
\]
由矩阵乘法，
\[
C_{ij}=\sum_{k=1}^n H_{ik}R_{kj}.
\]
若 \(R_{kj}\neq 0\)，则必须有
\[
k\leqslant j.
\]
因此
\[
C_{ij}=\sum_{k=1}^{j} H_{ik}R_{kj}.
\]
若 \(i>j+1\)，则对任意 \(k\leqslant j\)，都有
\[
i>j+1>j\geqslant k,
\]
即
\[
i>k+1.
\]
由于 \(H\) 是上 Hessenberg 矩阵，
\[
H_{ik}=0,\qquad i>k+1.
\]
所以对每个 \(k\leqslant j\)，都有 \(H_{ik}=0\)，从而
\[
C_{ij}=0.
\]
因此
\[
(HR)_{ij}=0,\qquad i>j+1.
\]
所以 \(HR\) 是上 Hessenberg 矩阵。
\[
\boxed{HR\text{ 是上 Hessenberg 矩阵}.}
\]
\end{solution}

\item 设 \(A\in\C^{n\times n}\)，\(x\in\C^n\)，
\[
X=[x,Ax,\dots,A^{n-1}x].
\]
证明：如果 \(X\) 非奇异，则 \(X^{-1}AX\) 是上 Hessenberg 矩阵。

\begin{solution}
记
\[
X=[x_1,x_2,\dots,x_n],
\]
其中
\[
x_k=A^{k-1}x,\qquad k=1,2,\dots,n.
\]
因为 \(X\) 非奇异，所以 \(\{x_1,\dots,x_n\}\) 构成 \(\C^n\) 的一组基。

计算
\[
AX=[Ax,A^2x,\dots,A^nx].
\]
由 Cayley--Hamilton 定理，\(A^n\) 可表示为
\[
A^nx=c_1x+c_2Ax+\cdots+c_nA^{n-1}x.
\]
因此
\[
AX=[x_2,x_3,\dots,x_n,c_1x_1+\cdots+c_nx_n].
\]
于是
\[
X^{-1}AX
=
X^{-1}[x_2,x_3,\dots,x_n,c_1x_1+\cdots+c_nx_n].
\]
因为
\[
X^{-1}x_k=e_k,\qquad k=1,2,\dots,n,
\]
所以
\[
X^{-1}AX
=
[e_2,e_3,\dots,e_n,c],
\]
其中
\[
c=(c_1,c_2,\dots,c_n)^{\top}.
\]
因此
\[
X^{-1}AX=
\begin{pmatrix}
0&0&\cdots&0&c_1\\
1&0&\cdots&0&c_2\\
0&1&\cdots&0&c_3\\
\vdots&\vdots&\ddots&\vdots&\vdots\\
0&0&\cdots&1&c_n
\end{pmatrix}.
\]
这个矩阵只有次对角线及其以上可能非零，因此是上 Hessenberg 矩阵。
\[
\boxed{X^{-1}AX\text{ 是上 Hessenberg 矩阵}.}
\]
\end{solution}

\item 设 \(H\in\C^{n\times n}\) 是可对角化的不可约上 Hessenberg 矩阵。证明：\(H\) 没有重特征值。

\begin{solution}
反证。设 \(H\) 有重特征值 \(\lambda\)，重数至少为 \(2\)。

因为 \(H\) 可对角化，所以对应于 \(\lambda\) 的特征子空间维数至少为 \(2\)。设
\[
Hx=\lambda x,\qquad Hy=\lambda y,
\]
其中 \(x,y\) 线性无关。令
\[
z=\alpha x+\beta y\neq 0,
\]
则
\[
Hz=\lambda z.
\]
因此对应于 \(\lambda\) 的特征子空间至少是二维的。

另一方面，对于不可约上 Hessenberg 矩阵，若 \(\lambda\) 是特征值，则其几何重数为 \(1\)。这是因为不可约上 Hessenberg 矩阵的次对角元均非零。设
\[
(H-\lambda I)v=0.
\]
若 \(v\neq 0\)，取 \(v\) 的第一个非零分量 \(v_k\)。由第 \(k-1\) 个方程
\[
h_{k-1,k}v_k+h_{k-1,k-1}v_{k-1}+\cdots=0
\]
以及前面的零分量可推出 \(v_{k-1}=0\)，矛盾。因此特征子空间维数为 \(1\)。

于是对应于 \(\lambda\) 的特征子空间维数只能为 \(1\)，与可对角化且有重特征值矛盾。

因此 \(H\) 没有重特征值。
\[
\boxed{H\text{ 没有重特征值}.}
\]
\end{solution}

\item 设 \(H\in\R^{n\times n}\) 是一个不可约的上 Hessenberg 矩阵。证明：存在一个对角阵
\(D\in\R^{n\times n}\)，使得 \(D^{-1}HD\) 的次对角元均为 \(1\)，并求
\[
\kappa_2(D)=\norm{D}_2\norm{D^{-1}}_2.
\]

\begin{solution}
设
\[
H=
\begin{pmatrix}
h_{11}&h_{12}&\cdots&h_{1n}\\
h_{21}&h_{22}&\cdots&h_{2n}\\
0&h_{32}&\cdots&h_{3n}\\
\vdots&\ddots&\ddots&\vdots\\
0&\cdots&0&h_{n,n-1}&h_{nn}
\end{pmatrix}.
\]
不可约意味着
\[
h_{i+1,i}\neq 0,\qquad i=1,2,\dots,n-1.
\]
取对角矩阵
\[
D=\diag(d_1,d_2,\dots,d_n),
\]
其中 \(d_1=1\)，并递推取
\[
d_{i+1}=\frac{d_i}{h_{i+1,i}},\qquad i=1,2,\dots,n-1.
\]
即
\[
d_1=1,\qquad
d_2=\frac{1}{h_{21}},\qquad
d_3=\frac{1}{h_{21}h_{32}},\qquad
\dots,\qquad
d_n=\frac{1}{h_{21}h_{32}\cdots h_{n,n-1}}.
\]
计算 \(D^{-1}HD\) 的次对角元。因为
\[
(D^{-1}HD)_{i+1,i}
=
d_{i+1}^{-1}h_{i+1,i}d_i
=
\frac{d_i}{d_{i+1}}h_{i+1,i}.
\]
由 \(d_{i+1}=d_i/h_{i+1,i}\)，得
\[
\frac{d_i}{d_{i+1}}=h_{i+1,i}.
\]
因此
\[
(D^{-1}HD)_{i+1,i}=h_{i+1,i}h_{i+1,i}^{-1}=1.
\]
所以 \(D^{-1}HD\) 的次对角元均为 \(1\)。

下面求 \(\kappa_2(D)\)。因为 \(D\) 是对角矩阵，
\[
\norm{D}_2=\max_{1\leqslant i\leqslant n}\abs{d_i},
\qquad
\norm{D^{-1}}_2=\max_{1\leqslant i\leqslant n}\abs{d_i}^{-1}.
\]
因此
\[
\kappa_2(D)=\norm{D}_2\norm{D^{-1}}_2
=
\left(\max_i\abs{d_i}\right)
\left(\max_i\abs{d_i}^{-1}\right).
\]
又由递推
\[
d_i=\frac{1}{h_{21}h_{32}\cdots h_{i,i-1}},\qquad i=1,\dots,n,
\]
其中空乘积为 \(1\)。因此
\[
\boxed{
\kappa_2(D)
=
\left(\max_{1\leqslant i\leqslant n}\abs{d_i}\right)
\left(\max_{1\leqslant i\leqslant n}\abs{d_i}^{-1}\right),
\qquad
d_i=\frac{1}{\prod_{k=1}^{i-1}h_{k+1,k}}.
}
\]
\end{solution}

\item 设 \(H\in\R^{n\times n}\) 是一个上 Hessenberg 矩阵，并假定
\(x\in\R^n\) 是 \(H\) 的对应于实特征值 \(\lambda\) 的一个特征向量。试给出一个算法计算正交矩阵
\(Q\in\R^{n\times n}\)，使得
\[
Q^{\top}HQ=
\begin{pmatrix}
\lambda&w^{\top}\\
0&H_1
\end{pmatrix},
\]
其中 \(H_1\) 是 \((n-1)\) 阶上 Hessenberg 矩阵。

\begin{solution}
算法如下。

\begin{enumerate}
\item 取单位特征向量 \(q_1\)，即
\[
q_1=\frac{x}{\norm{x}_2},
\]
使得
\[
Hq_1=\lambda q_1.
\]

\item 构造正交矩阵 \(Q_0\in\R^{n\times n}\)，使其第一列为 \(q_1\)。例如可取
\[
Q_0=[q_1,q_2,\dots,q_n],
\]
其中 \(q_2,\dots,q_n\) 是 \(q_1\) 的正交补的一组标准正交基。

\item 计算
\[
Q_0^{\top}HQ_0.
\]
由于 \(Hq_1=\lambda q_1\)，有
\[
Q_0^{\top}HQ_0
=
\begin{pmatrix}
\lambda&b^{\top}\\
0&B
\end{pmatrix},
\]
其中 \(B\in\R^{(n-1)\times(n-1)}\)。

\item 一般 \(B\) 不一定是上 Hessenberg 矩阵。对 \(B\) 做 Householder 相似变换，将其化为上 Hessenberg 矩阵。即求正交矩阵
\(V\in\R^{(n-1)\times(n-1)}\)，使得
\[
V^{\top}BV=H_1
\]
为上 Hessenberg 矩阵。

\item 令
\[
Q=Q_0
\begin{pmatrix}
1&0\\
0&V
\end{pmatrix}.
\]
则 \(Q\) 是正交矩阵，并且
\[
Q^{\top}HQ
=
\begin{pmatrix}
1&0\\
0&V^{\top}
\end{pmatrix}
Q_0^{\top}HQ_0
\begin{pmatrix}
1&0\\
0&V
\end{pmatrix}
=
\begin{pmatrix}
\lambda&w^{\top}\\
0&H_1
\end{pmatrix},
\]
其中
\[
w^{\top}=b^{\top}V,
\]
且 \(H_1\) 是 \((n-1)\) 阶上 Hessenberg 矩阵。
\end{enumerate}

因此算法输出为
\[
\boxed{
Q=Q_0
\begin{pmatrix}
1&0\\
0&V
\end{pmatrix},
\qquad
Q^{\top}HQ=
\begin{pmatrix}
\lambda&w^{\top}\\
0&H_1
\end{pmatrix}.
}
\]
\end{solution}

\item 设 \(H\) 是一个奇异的不可约上 Hessenberg 矩阵。证明：进行一次基本的 QR 迭代后，\(H\) 的零特征值将出现。

\begin{solution}
设 \(H\) 奇异，则
\[
\det(H)=0,
\]
所以 \(0\) 是 \(H\) 的特征值。

对 \(H\) 进行基本的 QR 迭代。先做 QR 分解：
\[
H=QR,
\]
其中 \(Q\) 是酉矩阵，\(R\) 是上三角矩阵。然后令
\[
H_1=RQ.
\]
由于
\[
H_1=RQ=Q^{H}HQ,
\]
所以 \(H_1\) 与 \(H\) 酉相似，因而有相同特征值，特别地 \(0\) 仍是 \(H_1\) 的特征值。

因为 \(H\) 是不可约上 Hessenberg 矩阵，所以其次对角元均非零。QR 迭代保持上 Hessenberg 形式，因此 \(H_1\) 也是上 Hessenberg 矩阵。

又因为 \(H\) 奇异，所以 \(R\) 奇异。由于 \(R\) 是上三角矩阵，故
\[
r_{nn}=0.
\]
于是
\[
H_1=RQ
\]
的最后一行为
\[
(H_1)_{n,:}=(RQ)_{n,:}=r_{nn}Q_{n,:}=0.
\]
因此 \(H_1\) 的最后一行全为零。于是
\[
H_1=
\begin{pmatrix}
*&*&\cdots&*&*\\
*&*&\cdots&*&*\\
\vdots&\vdots&\ddots&\vdots&\vdots\\
0&0&\cdots&*&*\\
0&0&\cdots&0&0
\end{pmatrix}.
\]
这样的矩阵具有零特征值 \(0\)，并且零特征值在 \(H_1\) 中已经“出现”在最后一行。

因此进行一次基本 QR 迭代后，\(H\) 的零特征值将出现。
\[
\boxed{\text{一次基本 QR 迭代后，零特征值出现}.}
\]
\end{solution}

\end{enumerate}

\end{document}